\section{Ejemplo de uso y explicación}

A continuación, se presenta un ejemplo detallado de cómo utilizar el solver para la optimización de funciones matemáticas o problemas binarios utilizando el archivo de configuración centralizado.

\subsection{Optimización de funciones matemáticas (Benchmark) y Problemas Binarios}

Para comenzar con la optimización, lo primero que deben hacer es configurar las instancias, algoritmos y parámetros en el archivo de configuración.

\subsubsection{Paso 1: Configuración del archivo \texttt{experiments\_config.json}}

Diríjanse al directorio \texttt{src/util/json/} y abran el archivo \texttt{experiments\_config.json}. Dentro de este archivo encontrarán la estructura que controla la ejecución. A continuación, se detallan las secciones clave:

\begin{itemize}

\item \textbf{Selección del tipo de problema}
\\
Para habilitar o deshabilitar problemas específicos, se utilizan llaves booleanas. Por ejemplo, para ejecutar optimización en Benchmark (BEN) y desactivar SCP, USCP y KP:
\begin{quote}
\begin{verbatim}
"ben": true, 
"scp": false, 
"uscp": false,
"kp": false,
\end{verbatim}
\end{quote}

\item \textbf{Selección de metaheurísticas}
\\
En la lista \texttt{"mhs"} se declaran las siglas de los algoritmos a ejecutar (por ejemplo, SBOA, PSO, GWO):
\begin{quote}
\begin{verbatim}
"mhs": ["SBOA"],
\end{verbatim}
\end{quote}

\item \textbf{Configuración de acciones de discretización}
\\
La lista \texttt{"DS\_actions"} define las binarizaciones para problemas discretos (como SCP, USCP o KP). Esto no influye en las funciones continuas (BEN), pero es crucial para los problemas discretos:
\begin{quote}
\begin{verbatim}
"DS_actions": ["V3-ELIT"],
\end{verbatim}
\end{quote}

\item \textbf{Configuración de Mapas Caóticos (Novedad)}
\\
El framework permite ejecutar experimentos usando mapas caóticos en las ecuaciones de perturbación. Se puede activar globalmente y definir modos como \texttt{"comparacion"} (ejecuta tanto el método estándar como el caótico para comparar resultados):
\begin{quote}
\begin{verbatim}
"usar_mapas_caoticos": true,
"mapas_caoticos": {
    "modo": "comparacion",
    "mapas_activos": ["LOG", "SINE"],
    "SCP": {
        "usar_caoticos": true,
        "modo": "comparacion",
        "mapas": ["LOG"]
    }
}
\end{verbatim}
\end{quote}

\item \textbf{Criterios de Parada y Parámetros de Ejecución}
\\
En la sección \texttt{"experimentos"} se especifica el criterio de parada global (\texttt{modo\_terminacion}: \texttt{"iter"} para iteraciones, \texttt{"fe"} para evaluaciones de función, o \texttt{"both"} para ambas), el límite global de evaluaciones (\texttt{"max\_fe"}) y los parámetros específicos para cada dominio:
\begin{quote}
\begin{verbatim}
"experimentos": {
    "modo_terminacion": "iter",
    "max_fe": 50000,
    "BEN": {
        "iteraciones": 500,
        "poblacion": 50,
        "num_experimentos": 1
    },
    "SCP": {
        "iteraciones": 100,
        "poblacion": 10,
        "num_experimentos": 31
    }
}
\end{verbatim}
\end{quote}

\item \textbf{Selección de instancias a optimizar}
\\
En la sección \texttt{"instancias"} se indica la lista exacta de instancias o funciones por cada tipo de problema. Por ejemplo, para Benchmark y SCP:
\begin{quote}
\begin{verbatim}
"instancias": {
    "BEN": ["F1CEC2017", "F2CEC2017", "F3CEC2017"],
    "SCP": ["41"],
    "USCP": ["u41"],
    "KP": []
}
\end{verbatim}
\end{quote}

\end{itemize}

\subsubsection{Paso 2: Poblar la base de datos de experimentos}

Una vez configurado el archivo \texttt{experiments\_config.json}, es necesario inyectar los experimentos parametrizados en la base de datos de SQLite antes de poder ejecutarlos. Para ello, ejecuten en la consola:
\begin{verbatim}
python scripts/crearBD.py
python scripts/poblarDB.py
\end{verbatim}
El script \texttt{poblarDB.py} se encargará de crear los registros de cada corrida para cada metaheurística e instancia seleccionada en la base de datos.

\subsubsection{Paso 3: Ejecutar el solucionador}

Una vez poblada la base de datos, tienen tres alternativas para ejecutar el solver según la conveniencia de su máquina:

\begin{enumerate}
    \item \textbf{Ejecución secuencial}:
    \begin{verbatim}
    python main.py
    \end{verbatim}
    Ejecuta un experimento a la vez secuencialmente mostrando el progreso por consola.
    
    \item \textbf{Ejecución paralela local (Joblib)}:
    \begin{verbatim}
    python main_joblib.py --n-jobs 4
    \end{verbatim}
    Utiliza el framework de paralelismo Joblib para competir de forma segura por los experimentos en la base de datos, acelerando drásticamente el proceso.
    
    \item \textbf{Ejecución distribuida en consolas (Windows)}:
    \begin{verbatim}
    python levantarCMD.py
    \end{verbatim}
    Este script abre múltiples terminales CMD ejecutando \texttt{main.py} de manera concurrente para paralelizar el consumo de la base de datos.
\end{enumerate}
